% ED26-SCI-NIST2012-001: complete mathematical derivations.

\begin{annexitem}{annex:p15.factorial.estimability}{thm:p15-factorial-estimability}{Estimable functionals and indistinguishable parameter values}
\begin{proof}
Linearity of expectation gives
\(\mathbb E[a^{\mathsf T}Y]=a^{\mathsf T}X\beta\).
Equality to \(c^{\mathsf T}\beta\) for every \(\beta\) is exactly
\(X^{\mathsf T}a=c\), which is membership of \(c\) in the row space.
Every vector in that row space is orthogonal to \(\ker X\).
Conversely, decompose \(c=c_0+v\) into its row-space component and its
orthogonal complement. The latter belongs to \(\ker X\), and
\(c^{\mathsf T}v=\|v\|^2\). If \(c\) annihilates the kernel, then
\(v=0\), proving the converse.

The equality \(X\beta=X\beta'\) means
\(\beta-\beta'\in\ker X\). This proves equivalence with invariance
of the target on fitted-mean fibers. If the condition fails, take the
nonzero \(v\) above: \(\beta\) and \(\beta+v\) have identical
fitted means and targets differing by \(\|v\|^2\). With a fixed noise
law, their complete observation laws agree. No statistic of those
observations can distinguish the two parameter values. Finally, duplicating
rows changes neither their span nor the kernel, proving the replication
claim. It can change variance by adding observations but not identifiability
in this model.
\end{proof}
\end{annexitem}

\begin{annexitem}{annex:p15.factorial.two-by-two}{thm:p15-factorial-two-by-two}{The four-cell contrasts under both codings}
\begin{proof}
Each nonintercept column of \(H_2\) has two plus and two minus entries.
The product of any two distinct columns is another balanced nonintercept
column. Every column has squared norm four. Thus
\(H_2^{\mathsf T}H_2=4I_4\), giving the inverse transformation.
Substituting \(z=2b-1\), the difference between \(x=1\) and \(x=-1\)
is \(2\beta_A+2(2b-1)\beta_{AB}\). Averaging over \(b\) cancels
the interaction term, and differencing the two conditional effects gives
\(4\beta_{AB}\). The same calculation with the factors exchanged proves
the formulas for B. Expanding \(x=2a-1,z=2b-1\) gives intercept
\(\mu_{00}\), linear coefficients
\(\mu_{10}-\mu_{00}\) and \(\mu_{01}-\mu_{00}\), and product
coefficient \(\mu_{11}-\mu_{10}-\mu_{01}+\mu_{00}\).

For the retained means, direct multiplication gives
\(H_2^{\mathsf T}(4,0,0,2)^{\mathsf T}/4
=(3/2,-1/2,-1/2,3/2)^{\mathsf T}\). The stated row weights then give
the aid-present value two, the equal-weight checking value minus one and
the difference of differences six. With cell 11 absent, the fourth
coordinate of the cell-mean vector is in the kernel of the observation
map. Every target with a nonzero weight on that coordinate therefore fails
the estimability criterion, independently of a choice of regression coding.
\end{proof}
\end{annexitem}

\begin{annexitem}{annex:p15.factorial.linear-uncertainty}{thm:p15-factorial-linear-uncertainty}{Least variance, residual degrees of freedom, and the Student pivot}
\begin{proof}
For any solution \(a\) of \(X^{\mathsf T}a=c\), let
\(a_*=P_Xa\). The orthogonal complement of \(\operatorname{col}(X)\)
is \(\ker X^{\mathsf T}\), so \(X^{\mathsf T}a_*=c\).
Two solutions in \(\operatorname{col}(X)\) differ by a vector in that
space and its orthogonal complement, hence coincide. Every solution is
\(a_*+v\) with \(v\perp\operatorname{col}(X)\). Consequently
\[
 \operatorname{Var}((a_*+v)^{\mathsf T}Y)
 =\sigma^2(\|a_*\|^2+\|v\|^2),
\]
which proves minimality and the variance formula. In the full-column-rank
case direct multiplication verifies the displayed expression for \(a_*\)
and the usual least-squares coefficients.

The residual vector is \((I-P_X)\epsilon\). For an orthogonal projection
of rank \(N-r\),
\[
 \mathbb E\|(I-P_X)\epsilon\|^2
 =\operatorname{tr}((I-P_X)\sigma^2I_N)=\sigma^2(N-r).
\]
This trace identity also follows by summing the variances of residual
coordinates in an orthonormal basis. It proves unbiasedness of \(s^2\)
without normality.

For Gaussian noise choose an orthonormal basis with its first \(r\)
vectors spanning \(\operatorname{col}(X)\), and with its first vector
equal to \(a_*/\|a_*\|\). The latter is defined because \(c\ne0\).
The density of \(\epsilon/\sigma\) is proportional to
\(\exp(-\|z\|^2/2)\); an orthogonal change of coordinates preserves
the norm and has Jacobian of absolute value one. Its coordinates are
therefore independent standard normals. The normalized contrast error is
the first coordinate, and \(\mathrm{SSE}/\sigma^2\) is the sum of
squares of the final \(N-r\) coordinates. They are independent, and
the latter has the \(\chi^2_{N-r}\) law. The Student law of
\cref{thm:p15-student-law} gives the pivot
\[
 \frac{a_*^{\mathsf T}Y-c^{\mathsf T}\beta}{s\|a_*\|}.
\]
Its central interval between the two stated quantiles has probability
\(1-\eta\), proving the coverage claim. The interval formula is defined
also at \(s=0\); that event has probability zero under these assumptions.
When \(N=r\), no residual coordinates remain and the construction does
not supply a variance estimate.
\end{proof}
\end{annexitem}

\begin{annexitem}{annex:p15.factorial.joint-test}{thm:p15-factorial-joint-test}{The nested-space Gaussian test}
\begin{proof}
For nested spaces, \(P_1-P_0\) is the orthogonal projection onto
\(V_1\cap V_0^\perp\), of dimension \(r_1-r_0\). Its range is
orthogonal to the range of \(I-P_1\), of dimension \(N-r_1\).
Under the null both projections annihilate the mean. Expressing the normal
error in an orthonormal basis adapted to these two spaces gives disjoint
sets of independent standard-normal coordinates. Their squared lengths,
divided by \(\sigma^2\), are independent chi-squared variables with
the stated dimensions. Their ratio after division by those dimensions is
the defined F law. The denominator is positive almost surely; define the
test arbitrarily on its null-probability zero event. The chi-squared laws
have positive continuous densities on the positive axis, so the ratio has
a continuous strictly increasing CDF there and the stated upper-tail
critical value exists. Its null rejection probability is \(\eta\).

For one-factor or complete-block additive models, the centered treatment
space has dimension \(L-1\) and is orthogonal to the retained intercept
and block spaces. Removing precisely that space imposes equal treatment
means and gives the stated degrees of freedom. A rejection addresses that
joint mean restriction; it does not choose a particular mechanism or say
that every pair of treatment levels differs.
\end{proof}
\end{annexitem}

\begin{annexitem}{annex:p15.factorial.replication}{thm:p15-factorial-replication}{Replication, pure error, and clustering}
\begin{proof}
The variance of an average of \(n_j\) independent common-variance
responses is \(\sigma^2/n_j\). Within a cell, subtracting its sample
mean is an orthogonal projection of rank \(n_j-1\); the preceding trace
calculation gives expectation \((n_j-1)\sigma^2\). Summing proves the
pure-error expectation.

Write each residual as
\[
 Y_{ji}-\widehat\mu_j
 =(Y_{ji}-\overline Y_j)+(\overline Y_j-\widehat\mu_j).
\]
The cross term sums to zero within the cell because
\(\sum_i(Y_{ji}-\overline Y_j)=0\). Squaring and summing gives the
decomposition. If a lower-rank fitted mean model is correct, the second
term has the remaining model lack-of-fit degrees of freedom; it cannot be
treated as pure error without the mean-model assumption.

For repeated measurements on one unit,
\(\overline Y_i=\mu+U_i+q^{-1}\sum_je_{ij}\).
Independence of components gives variance
\(\sigma_U^2+\sigma_e^2/q\). The unit term is shared by its repeated
measurements and is not divided by \(q\). This is exactly why treating
them as independently assigned units can understate uncertainty.
\end{proof}
\end{annexitem}

\begin{annexitem}{annex:p15.factorial.blocks}{thm:p15-factorial-blocks}{Additive complete-block estimation}
\begin{proof}
In the \(L\)-by-\(B\) grid, the intercept, the centered treatment
columns and the centered block columns form mutually orthogonal spaces.
For a treatment vector \(u_t\) with \(\sum_tu_t=0\), and a block
vector \(v_b\) with \(\sum_bv_b=0\), their inner product across the
grid is \((\sum_tu_t)(\sum_bv_b)=0\). The dimensions are
\(1,L-1,B-1\), giving rank \(L+B-1\).
Projection onto these spaces gives the grand mean, row means minus grand
mean, and column means minus grand mean. Alternatively, substituting these
estimates makes every residual row sum and column sum zero, exactly the
normal equations, and the centered parameterization is unique.

For \(\sum_tc_t=0\),
\[
 \sum_tc_t\overline Y_{t\cdot}
 =\sum_tc_t\tau_t+\frac1B\sum_{t,b}c_t\epsilon_{tb}.
\]
The intercept and the average of all block effects cancel. Independent
common-variance errors give variance
\(\sigma^2 B\sum_tc_t^2/B^2\). Subtracting the model rank from
\(LB\) gives \((L-1)(B-1)\). An unrestricted interaction has exactly
this residual subspace; with one observation per cell it can fit every
residual. Thus the additive-model residual estimate relies on absence of
those interactions and is not a separate pure-error estimate.

The one-factor result is the same projection with the block space omitted:
balanced group means project onto a grand mean and centered treatment
effects. The number of fixed-count randomization sequences follows by
permuting \(N\) distinguishable units and dividing by \(n_t!\) for
the interchangeable labels in each treatment group.
\end{proof}
\end{annexitem}

\begin{annexitem}{annex:p15.factorial.squares}{thm:p15-factorial-squares}{Orthogonal square factors and their information budget}
\begin{proof}
Every level of a factor occurs \(s\) times. A centered level function
therefore gives a contrast column orthogonal to the intercept. Its space
has dimension \(s-1\): a level function that vanishes on every run
vanishes at each level. For two factors \(j,\ell\), pairwise balance
gives inner product
\[
 \sum_{a=1}^{s}\sum_{b=1}^{s}\theta_j(a)\theta_\ell(b)
 =(\sum_a\theta_j(a))(\sum_b\theta_\ell(b))=0.
\]
The intercept and all \(k\) contrast spaces are thus a direct orthogonal
sum of dimension \(1+k(s-1)\). Projection onto a factor's contrast
space is its level mean minus the grand mean; contributions of the other
factors average to zero at each of its levels. Subtraction from \(s^2\)
gives the residual dimension. For \(s=3,k=5\) this would require rank
eleven in a nine-dimensional observation space, which is impossible.
The two stated saturated cases have ranks nine and sixteen, respectively.

The displayed prime-order squares can also be verified by modular
arithmetic. With zero-based row and column \(r,c\), a column
\(c+ar\) modulo \(s\) is balanced against either index for nonzero
\(a\). Two different slopes give every ordered pair once, since
their difference determines \(r\), and then \(c\). For the three-level
Graeco-Latin example the slopes are 1 and 2; for the five-level example
they are 1 and 2, and the Hyper-Graeco-Latin example uses 1, 2 and 3.
The Latin examples at levels three and five use slope 2. Adding one to
the residues and encoding the treatment by letters gives the displayed
tables. For the four-level Graeco-Latin and Hyper-Graeco-Latin tables,
use the field with four elements, encoded by
\(0,1,a,a+1\), where \(a^2=a+1\). The extra columns are distinct
nonzero slopes. Coefficients are calculated modulo two; reduction by
\(a^2+a+1\), which has no root in the two-element field, gives the
four-element field multiplication. The selected slopes are
1 and \(a+1\) for the Graeco-Latin example, and
\(a+1,1,a\) in the indicated order for the Hyper-Graeco-Latin example.
The same invertible two-equation calculation proves pairwise balance.
In the separately displayed four-level Latin square, each row and column
explicitly contains A, B, C and D once, which is the required balance
against its two index factors.
\end{proof}
\end{annexitem}

\begin{annexitem}{annex:p15.factorial.walsh}{thm:p15-factorial-walsh}{Product-column contrasts and the replicated source calculations}
\begin{proof}
The product \(\chi_S\chi_T\) equals \(\chi_{S\triangle T}\), since
every repeated factor has square one. For a nonempty set of factors,
flipping one factor in that set pairs each positive product value with
a negative value over the full cube. Its sum is zero. The empty product
sums to \(2^k\). This proves orthogonality and the coefficient formula.
Each nonempty product has exactly half its entries of each sign; hence
the difference of its two half-cube means is twice the coefficient.

For independent cell means of variance \(\sigma^2/n\), the variance
of their signed normalized sum is
\(2^{-2k}2^k\sigma^2/n=\sigma^2/(n2^k)\). Covariances of distinct
estimates vanish by the orthogonality sum. There are \(n-1\) pure-error
degrees of freedom at each of \(2^k\) cells.

For the source's two-replicate table the cell means are
\((-2,-1/2,-1/2,5/2,-1/2,3/2,1,11/2)\).
Multiplying by the displayed sign columns and dividing by eight gives
the eight coefficients stated in the body. For two observations in a
cell, the within-cell sum of squares is half their squared difference.
Their differences are \((-2,1,-1,-1,-1,1,0,1)\), whose squared sum
is ten. Thus pure-error sum of squares is five, with eight degrees of
freedom. The factorial-effect variance estimate is
\(4(5/8)/16=5/32\), as stated.

For the second source response vector, the signed sums in product-column
order are
\((414,34,6,-14,26,-58,2,14)\).
The intercept divides the first by eight; the effects divide the others
by four. On the negative half fraction, the sums for the three main
columns divide by two and yield \((8,16,10)\); on the positive half
they yield \((9,-13,3)\). The general alias theorem explains their
relation to the full factorial effects.
\end{proof}
\end{annexitem}

\begin{annexitem}{annex:p15.factorial.centers-blocks}{thm:p15-factorial-centers-blocks}{Block aliases and the information from a common center}
\begin{proof}
Multiplication of the three sign columns gives minus on standard rows
1, 4, 6 and 7 and plus on rows 2, 3, 5 and 8. Adding a block column with
those signs duplicates \(\chi_{123}\). The fitted contribution is
\((\gamma+\beta_{123})\chi_{123}\); changing the two coefficients
by opposite amounts leaves the mean unchanged. Every lower-order column
is orthogonal to that column in the full cube, proving the remaining
separation statements under the specified additive block model.
If the same center is run in both blocks, its two model means are
\(\mu_0-\gamma\) and \(\mu_0+\gamma\). Their difference identifies
\(2\gamma\), and subtracting it from the corner product-column contrast
identifies \(2\beta_{123}\). This cancellation permits a common center
curvature term but still assumes the stated additive block effect.

At the common center, linear, cross-product and square terms all vanish,
leaving \(\beta_0\). Across the full cube the linear and cross-product
terms average to zero, whereas every square equals one. The corner mean
is therefore \(\beta_0+\sum_jq_j\), proving the curvature contrast.
Each square column has value one on every corner and zero on the center,
so their differences lie in the analysis matrix's kernel. Individual
quadratic coefficients remain nonestimable. On the half fraction
\(x_3=x_1x_2\), both columns additionally have value zero at a common
center. They remain identical after any number of center replications.
This argument uses the actual enlarged matrix; the identity word
\(I=123\), valid at the corners, is not falsely imposed at a center.
\end{proof}
\end{annexitem}

\begin{annexitem}{annex:p15.factorial.alias}{thm:p15-factorial-alias}{Regular fractions as affine binary systems}
\begin{proof}
Write \(x_j=(-1)^{z_j}\) with \(z_j\) in the two-element field.
A generator becomes a binary linear equation
\(g_j\cdot z=b_j\), where \(s_j=(-1)^{b_j}\).
Independence of the generator words means that the matrix of these rows
has rank \(p\). Row reduction gives an affine solution space of
dimension \(k-p\), hence \(2^{k-p}\) settings. Products of the
generators correspond to their binary row span, giving \(2^p\) words
and a well-defined product sign for each.

If \(H\) is in that row span, \(\chi_H\) has constant value \(s(H)\)
on the affine solution space. If it is not, there exists a homogeneous
solution \(v\) to all the generator equations with \(H\cdot v=1\).
Indeed, otherwise \(H\) would annihilate their common null space and
belong to their row span by finite-dimensional row reduction. Translation
by \(v\) pairs the affine solutions and reverses \(\chi_H\), so its
sum is zero. Applying this with \(H=S\triangle T\) proves the displayed
inner-product formula. Constant ratio gives identical or opposite columns;
otherwise their inner product is zero. Expanding the full mean function
in these product columns and applying linearity gives the alias-sum formula.

For the principal three-factor fraction, the defining group has words
\(\varnothing,123\). Multiplying each word by 1, 2 and 3 gives exactly
the four listed alias classes. Its negative counterpart gives the same
classes with reversed signs. In particular, adding
\(t(e_3-e_{12})\) to the full coefficient vector changes neither the
principal-fraction fitted values nor its observation law with fixed noise.

For the first eight-factor example, multiplying
\(3456,12457,12358\) in all seven nonempty combinations gives precisely
the listed defining words. Their length-four members are
\(3456,3478,5678\). Splitting these into pairs produces
\(34=56=78\) and the six other paired classes in the body; all pairs
involving 1 or 2 are outside them. For the alternative, its five
length-four words are \(1357,2458,2467,1368,5678\).
Their pair splits give the eleven displayed classes, comprising
twenty-four two-factor interactions; the complement is exactly
\(12,14,23,34\). In particular \(27=46\) follows directly from
\(I=2467\).

Under the stated factor permutation, the first design's three generators
map to \(1236,23458,13578\). For the third displayed specification,
the generators are \(1236,1247,23458\), and
\(13578=1247\mathbin{\triangle}23458\). Each set therefore generates
the same binary word space, proving the claimed equivalence. All of these
are direct defining-group calculations, without assuming any effect is small.
\end{proof}
\end{annexitem}

\begin{annexitem}{annex:p15.factorial.resolution}{thm:p15-factorial-resolution}{What the shortest defining word excludes}
\begin{proof}
Two columns can be aliased only when their symmetric difference is a
nonempty defining word. Distinct singleton words differ in a word of size
two. A singleton and a two-factor word differ in size at most three;
two distinct two-factor words differ in size at most four. Thus resolutions
III, IV and V respectively exclude the aliasing relations claimed.
The existence assertions follow by splitting a shortest word: a word
of size three splits as one plus two factors, a word of size four as
two plus two, and a word of size five as two plus three. The resulting
two disjoint words have that defining word as their symmetric difference
and hence are aliased. The sign of the generator affects the sign of an
alias, not whether it exists.

At resolution at least V, any two distinct words of size at most two,
including the empty word, differ in fewer than five positions. They
therefore cannot belong to one alias class. The corresponding columns
are nonzero and mutually orthogonal, so the observation space must have
dimension at least their number, \(1+k+\binom{k}{2}\). With eight
factors the number is thirty-seven, and the least permissible power of
two is sixty-four. This is a rank bound, not an assumption that higher
interactions are absent.
\end{proof}
\end{annexitem}

\begin{annexitem}{annex:p15.factorial.assignment}{thm:p15-factorial-assignment}{Assignment expectations and their exact variance}
\begin{proof}
Potential responses and weights are fixed under this randomization.
Since \(\mathbb E[\mathbf1_{\{Z_i=t\}}]=\pi_i(t)\), taking expectation
of each term in \(\widehat\tau\) gives \(w_tY_i(t)/N\). Summing
proves unbiasedness. In complete randomization with cell sizes \(n_t\),
\(\pi_i(t)=n_t/N\), so the terms for cell \(t\) sum to
\(w_t n_t^{-1}\sum_{i:Z_i=t}Y_i(t)=w_t\overline Y_t\).

The covariance of two assignment indicators is
\(\pi_{ij}(t,u)-\pi_i(t)\pi_j(u)\). Expanding the variance of the
finite weighted sum proves the formula. This includes the cases
\(i=j,t=u\) and \(i=j,t\ne u\), where the joint probabilities are
\(\pi_i(t)\) and zero, respectively. Products of two potential outcomes
for the same unit under different assignments are not both observed;
the variance identity therefore does not by itself identify a data-based
variance estimator. Additional design bounds or an explicit outcome model
are needed for that step.
\end{proof}
\end{annexitem}

\begin{annexitem}{annex:p15.factorial.mechanism}{thm:p15-factorial-mechanism}{Two mechanisms with the same assignment surface}
\begin{proof}
Let \(A,B\) be assigned binary variables, \(R\) a putative reflection
variable and \(E\) a competing effort variable. In model one set
\[
 R=A,\quad E=0,\quad Y=4-4R-4B+6RB.
\]
In model two set
\[
 R=0,\quad E=A,\quad Y=4-4E-4B+6EB.
\]
These are finite acyclic structural systems, with outcomes among
\(0,2,4\). Under every simultaneous assignment intervention on \(A,B\),
both give \(Y=4-4A-4B+6AB\). They therefore agree on every factorial
assignment mean and all its contrasts, under any common randomization law.
But fixing \(A=0,B=0\) and additionally replacing the equation for
\(R\) by \(R=1\) gives \(Y=0\) in model one and \(Y=4\) in model
two. The variable targeted as reflection has different intervention effects.
The assignment response surface cannot distinguish the two models, even
if measured exactly. This is a direct nonidentification example under the
definition in \cref{def:p16-causal-identification}.
\end{proof}
\end{annexitem}
